% ============================================================
%  CHƯƠNG 6: KẾT LUẬN VÀ HƯỚNG PHÁT TRIỂN
% ============================================================
\chapter{Kết luận và hướng phát triển}
\label{chap:conclusion}

\section{Kết luận}
\label{sec:conclusion}

Tiểu luận đã trình bày toàn bộ quá trình thiết kế và hiện thực hệ thống khóa thông minh nhận diện khuôn mặt trên nền tảng Dual-MCU (STM32F411CEU6 + ESP32-S3). Các đóng góp chính của nhóm bao gồm:

\begin{enumerate}[label=\arabic*.]
  \item \textbf{Kiến trúc Dual-MCU hợp lý}: Phân tách rõ ràng tầng xử lý AI (ESP32-S3) và tầng điều khiển ngoại vi (STM32F411CEU6) giúp mỗi vi điều khiển tập trung vào thế mạnh của mình, đồng thời tăng tính mô-đun và khả năng bảo trì của hệ thống.
  \item \textbf{Giao thức UART nội bộ có tự phục hồi}: Thiết kế UART hai lớp giữa hai MCU giữ được khả năng debug ở bước bootstrap, đồng thời bổ sung xác thực gói, heartbeat và supervisor để tăng độ tin cậy vận hành thực tế.
  \item \textbf{Trải nghiệm người dùng đa kênh}: Tích hợp phản hồi qua OLED, âm thanh và LED tạo ra trải nghiệm rõ ràng, trực quan cho người dùng ở mọi điều kiện môi trường.
  \item \textbf{Vận hành độc lập, không phụ thuộc Internet/WAN}: Toàn bộ thuật toán nhận diện và cơ sở dữ liệu khuôn mặt được lưu trữ và xử lý trực tiếp trên thiết bị, đảm bảo tính riêng tư và hoạt động liên tục kể cả khi không có kết nối ra internet công cộng.
  \item \textbf{Bổ sung kênh điều khiển cục bộ qua Wi-Fi}: ESP32-S3 đồng thời cung cấp Web UI để preview camera, theo dõi trạng thái và phát lệnh mở khóa từ xa trong mạng nội bộ mà không cần thêm phần cứng mạng riêng.
\end{enumerate}

Kết quả thực nghiệm cho thấy hệ thống hoạt động ổn định, đạt tỉ lệ nhận diện chính xác cao trong điều kiện ánh sáng bình thường và thời gian phản hồi chấp nhận được cho ứng dụng kiểm soát truy cập thực tế.

\section{Hạn chế}
\label{sec:limitations}

Bên cạnh những kết quả đạt được, hệ thống hiện tại còn một số hạn chế cần được cải thiện:

\begin{enumerate}[label=\arabic*.]
  \item \textbf{Chỉ hỗ trợ 1 tư thế (chính diện)}: Hiện tại \texttt{ENROLL\_TOTAL\_STEPS = 1}, chỉ đăng ký khuôn mặt nhìn thẳng. Điều này giảm độ chính xác khi người dùng quay đầu hoặc thay đổi góc mặt so với đúc nhất đăng ký. Có thể nâng cấp lên 5 tư thế nhưng tốn thời gian đăng ký hơn.
  
  \item \textbf{Hiệu năng giảm trong điều kiện ánh sáng yếu}: Camera OV3660 không có khả năng tự động điều chỉnh ISO cao, dẫn đến ảnh nhiễu và tỉ lệ nhận diện giảm đáng kể khi độ chiếu sáng dưới 100 lux. Mặc dù hệ thống có đèn trợ sáng (GPIO47), nhưng độ sáng vẫn còn hạn chế so với camera với IR illuminator.
  
  \item \textbf{Không có cơ chế chống giả mạo (anti-spoofing)}: Hệ thống hiện tại có thể bị đánh lừa bằng ảnh in hoặc video phát lại trên điện thoại, vì chưa triển khai kỹ thuật phát hiện ``liveness'' dựa trên blink detection hoặc depth sensing.
  
  \item \textbf{Voting quá tối giản (\texttt{REQUIRED\_MATCHES = 1})}: Chỉ cần 1 frame khớp là mở cửa ngay, dễ xảy ra False Positive với ánh sáng phức tạp. Nếu tăng voting lên 2--3 frames, độ tin cậy tăng nhưng latency cũng tăng.
  
  \item \textbf{Bảo vệ UART mới dừng ở mức nhẹ}: Firmware hiện đã có bắt tay, sequence counter, auth tag và cơ chế che payload trên UART. Tuy vậy, khóa bí mật vẫn nằm trong flash firmware của hai MCU; nếu đối tượng tấn công có toàn quyền truy cập vật lý và trích xuất firmware, mức bảo vệ này chưa thể so sánh với secure element hoặc secure boot hoàn chỉnh.
  
  \item \textbf{Phục hồi ESP32 chưa xử lý được mọi nguyên nhân gốc}: Hiện tại STM32 đã có thể thử \texttt{STATUS}, gửi \texttt{REBOOT} và kéo reset cứng qua \texttt{PB1 -> EN/RESET}. Dù vậy, nếu nguyên nhân đến từ sụt áp nguồn 5~V, tiếp xúc GND kém hoặc lỗi camera phần cứng lặp lại, việc reset vẫn chỉ giải quyết triệu chứng chứ không loại bỏ hoàn toàn nguyên nhân nền.

  \item \textbf{Mở khóa từ xa qua Wi-Fi mới phù hợp cho mạng nội bộ tin cậy}: Web UI cục bộ hiện ưu tiên tính đơn giản để demo và vận hành trong lab. Endpoint \texttt{/unlock} chưa có lớp đăng nhập riêng, chưa dùng TLS, và còn tồn tại tùy chọn fallback AP mở nếu cấu hình bảo vệ thất bại. Vì vậy, tính năng này chưa đủ mạnh cho môi trường có đối tượng truy cập mạng không đáng tin cậy.
  
  \item \textbf{Dữ liệu khuôn mặt không được backup}: Cơ sở dữ liệu 7 khuôn mặt được lưu trong NVS Flash của ESP32-S3. Nếu flash bị lỗi hoặc bị xóa tình cờ, toàn bộ dữ liệu khuôn mặt sẽ mất mà không có cách phục hồi (không có backup on-cloud hay on-device ngoài).
\end{enumerate}

\section{Hướng phát triển}
\label{sec:future_work}

Từ những hạn chế trên, nhóm đề xuất các hướng phát triển tiếp theo theo mức ưu tiên:

\subsection{Ưu tiên cao (High Priority)}

\begin{enumerate}[label=\arabic*.]
  \item \textbf{Tăng voting thành 2--3 frames}: Thay đổi \texttt{REQUIRED\_MATCHES} từ 1 lên 2 hoặc 3 để giảm False Positive. Trade-off: latency tăng từ $\sim$130 ms lên 260--400 ms, không quá lớn.
  
  \item \textbf{Nâng cấp lên 5-pose enrollment}: Đặt \texttt{ENROLL\_TOTAL\_STEPS = 5} để yêu cầu 5 tư thế (FRONT, LEFT, RIGHT, UP, DOWN). Cải thiện đáng kể độ chính xác khi người dùng quay đầu, nhưng tăng thời gian đăng ký lên $\sim$50 giây.
  
  \item \textbf{Bổ sung đèn trợ sáng IR}: Lắp thêm IR LED (850 nm hoặc 940 nm) song song với visible LED (GPIO47) để cải thiện hiệu năng trong ánh sáng yếu. Cần điều chỉnh camera exposure settings.
  
  \item \textbf{Triển khai Liveness Detection}: Tích hợp mô hình nhỏ gọn dựa on blink detection hoặc optical flow để phát hiện giả mạo video/ảnh in. Có thể dùng MiniVGG hoặc Face Anti-Spoofing CVPR 2014 model được lượng tử hóa.
\end{enumerate}

\subsection{Ưu tiên trung (Medium Priority)}

\begin{enumerate}[label=\arabic*.]
  \item \textbf{Nâng cấp bảo vệ khoá bí mật UART}: Kết hợp Secure Boot / Flash Encryption của ESP32-S3 và RDP của STM32 để bảo vệ khóa bí mật hiện đang lưu trong firmware.
  
  \item \textbf{Gia cố xác thực cho Web UI mở khóa từ xa}: Bổ sung đăng nhập, token phiên, mật khẩu riêng theo thiết bị và vô hiệu hóa fallback AP mở trong bản triển khai thực tế.
  
  \item \textbf{Backup \& restore khuôn mặt}: Triển khai cơ chế xuất/nhập dữ liệu khuôn mặt qua UART hoặc USB, cho phép backup lên máy tính hoặc SD card bên ngoài.
  
  \item \textbf{Supervisor thông minh hơn}: Bổ sung ghi log nguyên nhân reset, đếm số lần phục hồi liên tiếp và phân loại lỗi brownout / camera / UART để tránh reset lặp vô ích.
  
  \item \textbf{Tích hợp NTP + Timestamping}: Sử dụng Wi-Fi của ESP32 để đồng bộ giờ từ NTP server, ghi timestamp vào log truy cập.
\end{enumerate}

\subsection{Ưu tiên thấp (Low Priority)}

\begin{enumerate}[label=\arabic*.]
  \item \textbf{Ứng dụng di động (Mobile App)}: Phát triển app iOS/Android để thay thế Web UI thử nghiệm bằng giao diện quản trị hoàn chỉnh hơn qua BLE hoặc Wi-Fi.
  
  \item \textbf{Nâng cấp RTOS cho STM32}: Chuyển từ superloop sang FreeRTOS để xử lý đa nhiệm thực sự, khử hoàn toàn blocking delay.
  
  \item \textbf{PCB thiết kế riêng}: Thay thế breadboard + dev boards bằng PCB để tăng độ tin cậy và giảm kích thước.
  
  \item \textbf{Machine Learning tự thích ứng}: Huấn luyện bổ sung embeddings trong quá trình vận hành để điều chỉnh theo thay đổi độ sáng hoặc góc độ mới của người dùng.
\end{enumerate}
